\documentclass[11pt,a4paper]{article}

\usepackage[a4paper,margin=1in]{geometry}
\usepackage{amsmath}
\usepackage{amssymb}
\usepackage{array}
\usepackage{booktabs}
\usepackage{enumitem}
\usepackage{tabularx}
\usepackage[T1]{fontenc}
\usepackage{lmodern}
\usepackage{xcolor}

\definecolor{packblue}{RGB}{26,68,114}
\definecolor{packgrey}{RGB}{85,85,85}

\newcommand{\markalloc}[1]{\nobreak\hfill\mbox{(#1 Marks)}}
\newcommand{\questiontotal}[1]{\begin{flushright}\textbf{Total #1 Marks}\end{flushright}}

\setlength{\parindent}{0pt}
\setlength{\parskip}{0.45em}
\setlist[enumerate,1]{label=(\alph*), leftmargin=1.1cm, itemsep=0.55em, topsep=0.4em}

\begin{document}

\begin{center}
{\Large\bfseries\color{packblue} Week 10 Revision Selected Questions}\\[0.4em]
{\large 5CM522 Operational Research}\\[0.3em]
{\color{packgrey} Reformatted from \texttt{Sem 2 Week 10 - Revision.pdf}}
\end{center}

\vspace{0.8em}
\hrule
\vspace{0.8em}

This document contains all selected questions from the Week 10 revision sheet, reformatted as a clean booklet.
The questions have been renumbered so that the pack reads continuously.

\newpage

\textbf{1.}

GreenGrid Energy is choosing a set of improvement projects for a regional depot network.
Each project is either undertaken in full or not undertaken at all.

\begin{center}
\begin{tabular}{lcc}
\toprule
Project & Cost (\pounds m) & Annual Benefit Score \\
\midrule
$A$ & 6 & 14 \\
$B$ & 8 & 17 \\
$C$ & 5 & 11 \\
$D$ & 9 & 20 \\
$E$ & 4 & 9 \\
$F$ & 7 & 15 \\
\bottomrule
\end{tabular}
\end{center}

The total budget is \pounds21 million.

The following additional conditions apply:
\begin{itemize}[leftmargin=0.95cm, itemsep=0.25em]
\item if project $D$ is selected, project $B$ must also be selected;
\item projects $A$ and $F$ cannot both be selected;
\item at least one of projects $C$ and $E$ must be selected.
\end{itemize}

\begin{enumerate}
\item Formulate a binary integer programming model to maximise total annual benefit subject to the budget and project-selection rules. Clearly define all decision variables. \markalloc{8}

\item Solve the model in Python or by systematic analysis, and determine which projects should be selected and the maximum total annual benefit score. \markalloc{7}

\item Explain why this problem requires binary decision variables rather than the continuous variables used in ordinary linear programming. \markalloc{3}

\item Modify the model so that exactly four projects must be selected, then resolve or explain how the optimal solution would change. \markalloc{4}

\item Suggest one realistic extension that would make this capital-selection model more useful in practice. No full formulation is required. \markalloc{3}
\end{enumerate}

\questiontotal{25}

\newpage

\textbf{2.}

Hock-Brooke (Automobiles), the main dealer in Ceeham, receives its new stock on Monday.
Most customers collect their new cars on Friday or Saturday.
New cars are stored initially in an old warehouse to be prepared for display in the showroom or, when the warehouse is full, in a secure parking area.

The dealership is about to replace its old large warehouse with a modern purpose-built warehouse and needs to know the most suitable size for this new warehouse.

The following data are available:

\begin{center}
\begin{tabular}{lcccc}
\toprule
Number of Cars Delivered & 7 & 8 & 9 & 10 \\
\midrule
Frequency (\%) & 10 & 20 & 50 & 20 \\
\bottomrule
\end{tabular}
\end{center}

\begin{center}
\begin{tabular}{lcccccc}
\toprule
Weekly Demand & 6 & 7 & 8 & 9 & 10 & 11 \\
\midrule
Frequency (\%) & 5 & 20 & 30 & 10 & 20 & 15 \\
\bottomrule
\end{tabular}
\end{center}

The cost of providing a single warehouse space is \pounds10 per week, and the cost of hiring an additional parking space is \pounds20 per week.

\begin{enumerate}
\item Simulate the operation of this dealership with a warehouse capacity of 9 cars and hence determine the average cost of using the warehouse per week. \markalloc{15}

\item Is 9 the best capacity? If not, determine the best possible capacity for the warehouse. \markalloc{10}
\end{enumerate}

\questiontotal{25}

\newpage

\textbf{3.}

A university library help desk is modelled as a queueing system.
Students arrive randomly and service times are assumed to follow an exponential distribution.
You should adapt the SimPy structures used in class, but tailor the model to this setting.

Use the following assumptions:

\begin{center}
\begin{tabular}{lc}
\toprule
Quantity & Value \\
\midrule
Simulation length & 360 minutes \\
Average service time & 7 minutes \\
Baseline arrival rate & 10 students per hour \\
Peak arrival rate & 15 students per hour \\
Service target & 97\% wait less than 8 minutes \\
Priority users & 10\% of arrivals \\
Suggested random seed & 24 \\
\bottomrule
\end{tabular}
\end{center}

\begin{enumerate}
\item Describe the main elements of a discrete-event simulation model for this system: entities, resources, events, stochastic inputs and performance measures. \markalloc{5}

\item Adapt a SimPy model to represent the baseline system with 2 advisers and an arrival rate of 10 students per hour. Report the total served, average wait, maximum wait and percentage waiting less than 8 minutes. \markalloc{6}

\item Investigate the peak-period scenario with 2, 3 and 4 advisers. Present your results in a clearly labelled table and determine the minimum number of advisers needed to meet the service target. \markalloc{4}

\item Extend the model so that 10\% of users are priority users who should be served ahead of standard users. Explain how the model should be changed and discuss the likely effect on standard-user waits. \markalloc{6}

\item Explain why a single run is not enough to support a staffing decision and describe how you would use replications and/or a warm-up period to improve the analysis. \markalloc{4}
\end{enumerate}

\questiontotal{25}

\newpage

\textbf{4.}

F.B. Knight (Deliveries) distributes goods from its warehouses to the WalShop supermarkets throughout the East Midlands.
To optimise profits for this fixed-price contract, the company needs to determine its best delivery schedule from warehouse to shop.

\begin{enumerate}
\item Use the following data to determine F.B. Knight's optimal distribution plan.

\begin{center}
\begin{tabular}{lccc}
\toprule
Warehouse & $A$ & $B$ & $C$ \\
\midrule
Available Stock & 100 & 100 & 100 \\
\bottomrule
\end{tabular}
\end{center}

\begin{center}
\begin{tabular}{lccc}
\toprule
Shop & $X$ & $Y$ & $Z$ \\
\midrule
Demand & 70 & 80 & 150 \\
\bottomrule
\end{tabular}
\end{center}

The shipping costs per 10 items are:

\begin{center}
\begin{tabular}{lccc}
\toprule
 & $X$ & $Y$ & $Z$ \\
\midrule
$A$ & 8 & 7 & 5 \\
$B$ & 12 & 3 & 9 \\
$C$ & 10 & 11 & 7 \\
\bottomrule
\end{tabular}
\end{center}

\markalloc{10}

\item Determine how the solution changes if the following additional restrictions are imposed:
\[
B \text{ cannot deliver to } Y
\]
and
\[
A \text{ must deliver at least 5 items to } Y.
\]
\markalloc{5}

\item Two new distribution locations act as intermediary nodes between the warehouses and shops and become available for use.
Determine whether you should use one or both locations.
Assume that if you use the locations, you do not ship directly from warehouses to shops.

Warehouse to intermediary costs per 10 items:

\begin{center}
\begin{tabular}{lcc}
\toprule
 & $D_1$ & $D_2$ \\
\midrule
$A$ & 4 & 5 \\
$B$ & 3 & 2 \\
$C$ & 6 & 4 \\
\bottomrule
\end{tabular}
\end{center}

Intermediary to shop costs per 10 items:

\begin{center}
\begin{tabular}{lcc}
\toprule
 & $D_1$ & $D_2$ \\
\midrule
$X$ & 4 & 2 \\
$Y$ & 5 & 4 \\
$Z$ & 2 & 4 \\
\bottomrule
\end{tabular}
\end{center}

\markalloc{10}
\end{enumerate}

\questiontotal{25}

\newpage

\textbf{5.}

The data analyst at N\&F has started an initial analysis of sales data within the firm.
The following table shows quarterly sales data over a 4-year period.

\begin{center}
\begin{tabular}{ccc}
\toprule
Period & Season & Sales Data \\
\midrule
1 & Spring & 1013 \\
2 & Summer & 1025 \\
3 & Autumn & 1002 \\
4 & Winter & 1021 \\
5 & Spring & 1053 \\
6 & Summer & 1067 \\
7 & Autumn & 1042 \\
8 & Winter & 1060 \\
9 & Spring & 1093 \\
10 & Summer & 1106 \\
11 & Autumn & 1083 \\
12 & Winter & 1101 \\
13 & Spring & 1134 \\
14 & Summer & 1147 \\
15 & Autumn & 1123 \\
16 & Winter & 1139 \\
\bottomrule
\end{tabular}
\end{center}

Use this data to compare and contrast the use of both additive and multiplicative models to forecast sales.

\begin{enumerate}
\item Plot the sales data and smoothed (trend) data. \markalloc{2}

\item Calculate seasonal variations using both additive and multiplicative models for seasonal variation. \markalloc{8}

\item Produce sales forecasts for next year using both models. \markalloc{12}

\item Evaluate the results from these forecasting models and make a recommendation concerning the most appropriate model to be used to forecast future sales. \markalloc{3}
\end{enumerate}

\questiontotal{25}

\vfill

\begin{center}
\textit{End of Selected Questions}
\end{center}

\end{document}
